%% kapa: exam=1 round=31 year=2020 period=1 form=A subject=civil_law subject_ko=민법 seq=27 number=27
%% source: Q-Net 감정평가사 기출문제 (cst00309) artlSeq=5207219 file=extracted/1008/5207219/A형/A형/2020년도 제31회 감정평가사 1차 시험_1교시_A형.pdf
%% conversion: kordoc 4.24.51 -> tools/gen1.py (auto); review=proofread-sonnet; flags=none
\begin{question}{27}
점유에 관한 설명으로 옳지 않은 것은? (다툼이 있으면 판례에 따름)
\begin{choices}
  \item 점유매개자의 점유는 자주점유이다.
  \item 점유는 사실상 지배로 성립한다.
  \item 다른 사정이 없으면, 건물의 소유자가 그 부지를 점유하는 것으로 보아야 한다.
  \item 점유매개관계가 소멸하면 간접점유자는 직접점유자에게 점유물의 반환을 청구할 수 있다.
  \item 점유자는 소유의 의사로 점유한 것으로 추정한다.
\end{choices}
\end{question}
\answer{1}
